\documentclass{../../../templates/natsirt}

\usepackage{pgfplots}
\pgfplotsset{compat=1.18}

% Course Metadata
\title{Calculus 2}
\author{Yi Hong}
\subtitle{Note:p}

% Layout and Geometry Adjustments
\setlength{\parindent}{0pt}
\setlength{\parskip}{3pt}
\setlength{\headheight}{15pt}

% Custom Box Environments for Pure Active-Recall Checklist
\newcounter{topic}[section]

\newenvironment{recallbox}[1][]{%
    \begin{bluebox}[#1]\refstepcounter{topic}\textbf{Definitions \& Concepts to State \thesection.\thetopic}%
}{\end{bluebox}}

\newenvironment{theoremsbox}[1][]{%
    \begin{greenbox}[#1]\refstepcounter{topic}\textbf{Theorems \& Identities to State \thesection.\thetopic}%
}{\end{greenbox}}

\newenvironment{challengebox}[1][]{%
    \begin{redbox}[#1]\refstepcounter{topic}\textbf{Proof Challenges \& Calculations \thesection.\thetopic}%
}{\end{redbox}}

\newenvironment{warningbox}[1][]{\begin{whitebox}[#1]\textbf{Conventions \& Exam Pitfalls. }}{\end{whitebox}}
\newenvironment{reviewbox}[1][]{\begin{whitebox}[#1]\textbf{Textbook Exercise Checklist. }}{\end{whitebox}}

% Checklist Item Macros
\newcommand{\citem}{\item[$\square$]}
\newcommand{\citemdone}{\item[$\text{\rlap{$\checkmark$}}\square$]}
\newcommand{\Img}{\operatorname{Im}}

\begin{document}

% =========================================================================
% PART 1: CHAPTER 11 — INFINITE SEQUENCES AND SERIES
% =========================================================================

% -------------------------------------------------------------------------
% SECTION 11.1: SEQUENCES
% -------------------------------------------------------------------------
\section{Sequences (Section 11.1)}

\begin{recallbox}
    \begin{itemize}
        \citem \textbf{Sequence Convergence:} State the formal $\epsilon$-$N$ definition of a sequence converging to limit $L$ ($\lim_{n \to \infty} a_n = L \iff \forall \epsilon > 0, \, \exists N \in \mathbb{N} \text{ such that } n > N \implies |a_n - L| < \epsilon$).
    \end{itemize}
\end{recallbox}

\begin{theoremsbox}
    \begin{itemize}
        \citem \textbf{Limit of a Function vs. Limit of a Sequence:} Notice that if $\lim_{x \to \infty} f(x) = L$ and $f(n) = a_n$ when $n$ is an integer, then:
        \[ \lim_{x \to \infty} f(x) = \lim_{n \to \infty} f(n) = L \]
        \citem \textbf{Limit Laws for Sequences:} Notice that the properties of limits (Sum, Difference, Constant Multiple, Product, Quotient, and Power Laws) are also satisfied for sequences.
        \citem \textbf{The Squeeze Theorem for Sequences:} State the Squeeze Theorem for sequences: if $a_n \le b_n \le c_n$ for all $n \ge n_0$ and $\lim_{n \to \infty} a_n = \lim_{n \to \infty} c_n = L$, then $\lim_{n \to \infty} b_n = L$.
    \end{itemize}
\end{theoremsbox}

\begin{challengebox}
    \begin{itemize}
        \citem \textbf{Absolute Value Limit Equivalence:} Prove that:
        \[ \lim_{n \to \infty} |a_n| = 0 \iff \lim_{n \to \infty} a_n = 0 \]
        using the Squeeze Theorem with $-|a_n| \le a_n \le |a_n|$.
    \end{itemize}
\end{challengebox}

\begin{theoremsbox}
    \begin{itemize}
        \citem \textbf{Continuous Function on a Sequence:} If $\lim_{n \to \infty} a_n = L$ and the function $f$ is continuous at $L$, then:
        \[ \lim_{n \to \infty} f(a_n) = f(L) \]
    \end{itemize}
\end{theoremsbox}

\begin{recallbox}
    \begin{itemize}
        \citem \textbf{Monotonic Sequences:} Write the definitions of a strictly monotonic sequence and a monotonic sequence (increasing, decreasing, strictly increasing, strictly decreasing).
        \citem \textbf{Bounded Sequences:} Write the definitions of a sequence bounded above ($\exists M$ such that $a_n \le M$), bounded below ($\exists m$ such that $a_n \ge m$), and the definition of a bounded sequence ($m \le a_n \le M$).
    \end{itemize}
\end{recallbox}

\begin{theoremsbox}
    \begin{itemize}
        \citem \textbf{Monotonic Sequence Theorem:} State the Monotonic Sequence Theorem: every bounded, monotonic sequence is convergent.
    \end{itemize}
\end{theoremsbox}

\begin{challengebox}
    \begin{itemize}
        \citem \textbf{Recursive Sequence Convergence and Limit Calculation:} Prove that the sequence:
        \[ a_1 = 2, \quad a_{n+1} = \frac{a_n + 6}{2}, \quad n \ge 1 \]
        is convergent, and find its limit $L$.
    \end{itemize}
\end{challengebox}

\begin{reviewbox}
    \begin{itemize}
        \citem \textbf{Textbook Exercise Revision (Section 11.1):} Problems 31, 37, 51.
        \begin{itemize}
            \citem Exercise 31
            \citem Exercise 37
            \citem Exercise 51
        \end{itemize}
    \end{itemize}
\end{reviewbox}

% -------------------------------------------------------------------------
% SECTION 11.2: SERIES
% -------------------------------------------------------------------------
\section{Series (Section 11.2)}

\begin{recallbox}
    \begin{itemize}
        \citem \textbf{Series:} What is an infinite series $\sum_{n=1}^\infty a_n$?
        \citem \textbf{Partial Sums and Convergence:} Write the definition of the partial sum $s_k = \sum_{n=1}^k a_n$, and define when a series converges to sum $s$ or diverges.
    \end{itemize}
\end{recallbox}

\begin{theoremsbox}
    \begin{itemize}
        \citem \textbf{Test for Divergence:} Write the Test for Divergence (if $\lim_{n \to \infty} a_n$ does not exist or $\lim_{n \to \infty} a_n \neq 0$, then the series $\sum_{n=1}^\infty a_n$ diverges), and prove it using $a_n = s_n - s_{n-1}$.
    \end{itemize}
\end{theoremsbox}

\begin{challengebox}
    \begin{itemize}
        \citem \textbf{Harmonic Series Divergence:} Show that the harmonic series $\sum_{n=1}^\infty \frac{1}{n}$ is divergent.
    \end{itemize}
\end{challengebox}

\begin{warningbox}
    \begin{itemize}
        \citem \textbf{Constant Multiple of Divergent Series:} Notice that if $\sum_{n=1}^\infty a_n$ diverges and $c \ne 0$, then $\sum_{n=1}^\infty c a_n$ also diverges.
    \end{itemize}
\end{warningbox}

\begin{reviewbox}
    \begin{itemize}
        \citem \textbf{Textbook Exercise Revision (Section 11.2):} Problems 15, 19, 31, 37, 39, 45, 49, 71.
        \begin{itemize}
            \citem Exercise 15
            \citem Exercise 19
            \citem Exercise 31
            \citem Exercise 37
            \citem Exercise 39
            \citem Exercise 45
            \citem Exercise 49
            \citem Exercise 71
        \end{itemize}
    \end{itemize}
\end{reviewbox}

% -------------------------------------------------------------------------
% SECTION 11.3: THE INTEGRAL TEST AND ESTIMATES OF SUMS
% -------------------------------------------------------------------------
\section{The Integral Test and Estimates of Sums (Section 11.3)}

\begin{theoremsbox}
    \begin{itemize}
        \citem \textbf{The Integral Test:} Write the Integral Test, and state all required conditions on the function $f$ (continuous, positive, and decreasing on $[1, \infty)$ where $f(n) = a_n$).
    \end{itemize}
\end{theoremsbox}

\begin{challengebox}
    \begin{itemize}
        \citem \textbf{Integral Test Application:} Determine whether:
        \[ \sum_{n=1}^\infty \frac{1}{n^2 + 1} \]
        is convergent or divergent.
    \end{itemize}
\end{challengebox}

\begin{recallbox}
    \begin{itemize}
        \citem \textbf{$p$-Series Definition:} Write the definition of a $p$-series $\sum_{n=1}^\infty \frac{1}{n^p}$.
    \end{itemize}
\end{recallbox}

\begin{challengebox}
    \begin{itemize}
        \citem \textbf{$p$-Series Convergence Proof:} Prove that the $p$-series is convergent $\iff p > 1$.
    \end{itemize}
\end{challengebox}

\begin{theoremsbox}
    \begin{itemize}
        \citem \textbf{Remainder and Sum Estimation:} If the series is convergent, we can approximate the sum $S$ by the remainder (error) $R_n$. Write the definition of $R_n = S - s_n$, and write the upper and lower bounds of $R_n$ and the estimate for $S$:
        \[ \int_{n+1}^\infty f(x)\,dx \le R_n \le \int_n^\infty f(x)\,dx \]
        \[ s_n + \int_{n+1}^\infty f(x)\,dx \le S \le s_n + \int_n^\infty f(x)\,dx \]
    \end{itemize}
\end{theoremsbox}

\begin{reviewbox}
    \begin{itemize}
        \citem \textbf{Textbook Exercise Revision (Section 11.3):} Problems 7, 9, 11, 17, 25, 29, 30, 31, 41.
        \begin{itemize}
            \citem Exercise 7
            \citem Exercise 9
            \citem Exercise 11
            \citem Exercise 17
            \citem Exercise 25
            \citem Exercise 29
            \citem Exercise 30
            \citem Exercise 31
            \citem Exercise 41
        \end{itemize}
    \end{itemize}
\end{reviewbox}

% -------------------------------------------------------------------------
% SECTION 11.4: THE COMPARISON TESTS
% -------------------------------------------------------------------------
\section{The Comparison Tests (Section 11.4)}

\begin{theoremsbox}
    \begin{itemize}
        \citem \textbf{Direct Comparison Test:} Write the Direct Comparison Test for series with positive terms ($a_n \le b_n \implies (\sum b_n \text{ converges} \implies \sum a_n \text{ converges})$ and $a_n \ge b_n \implies (\sum b_n \text{ diverges} \implies \sum a_n \text{ diverges})$).
    \end{itemize}
\end{theoremsbox}

\begin{challengebox}
    \begin{itemize}
        \citem \textbf{Direct Comparison Application:} Determine whether:
        \[ \sum_{n=1}^\infty \frac{\ln n}{n} \]
        is divergent or convergent by using the Direct Comparison Test.
        \citem \textbf{Limit Comparison Test and Proof:} Write the Limit Comparison Test (for series with positive terms $a_n, b_n > 0$, if $\lim_{n \to \infty} \frac{a_n}{b_n} = c \in (0, \infty)$), and prove it.
    \end{itemize}
\end{challengebox}

\begin{reviewbox}
    \begin{itemize}
        \citem \textbf{Textbook Exercise Revision (Section 11.4):} Problems 3, 5, 17, 21.
        \begin{itemize}
            \citem Exercise 3
            \citem Exercise 5
            \citem Exercise 17
            \citem Exercise 21
        \end{itemize}
    \end{itemize}
\end{reviewbox}

% -------------------------------------------------------------------------
% SECTION 11.5: ALTERNATING SERIES
% -------------------------------------------------------------------------
\section{Alternating Series (Section 11.5)}

\begin{recallbox}
    \begin{itemize}
        \citem \textbf{Alternating Series Definition:} Write the definition of an alternating series ($\sum (-1)^{n-1} b_n$ or $\sum (-1)^n b_n$ with $b_n > 0$).
    \end{itemize}
\end{recallbox}

\begin{theoremsbox}
    \begin{itemize}
        \citem \textbf{Alternating Series Test:} Write the Alternating Series Test (Leibniz's Test: if (i) $b_{n+1} \le b_n$ for all $n$, and (ii) $\lim_{n \to \infty} b_n = 0$, then the series converges).
    \end{itemize}
\end{theoremsbox}

\begin{warningbox}
    \begin{itemize}
        \citem \textbf{Alternating Series Error Bound:} Notice that for an alternating series satisfying the Alternating Series Test:
        \[ |s - s_n| = |R_n| \le b_{n+1} \]
    \end{itemize}
\end{warningbox}

\begin{recallbox}
    \begin{itemize}
        \citem \textbf{Absolute and Conditional Convergence:} Write the definitions of absolute convergence (when $\sum |a_n|$ converges) and conditional convergence (when $\sum a_n$ converges but $\sum |a_n|$ diverges).
    \end{itemize}
\end{recallbox}

\begin{challengebox}
    \begin{itemize}
        \citem \textbf{Absolute Convergence Test and Proof:} Write the Absolute Convergence Test (if $\sum |a_n|$ converges, then $\sum a_n$ converges), and prove it using $0 \le a_n + |a_n| \le 2|a_n|$ and the Direct Comparison Test.
    \end{itemize}
\end{challengebox}

\begin{theoremsbox}
    \begin{itemize}
        \citem \textbf{Rearrangement Theorem:} Write the Rearrangement Theorem (if $\sum a_n$ is an absolutely convergent series with sum $s$, then any rearrangement has the same sum $s$; if conditionally convergent, terms can be rearranged to sum to any real number or diverge).
    \end{itemize}
\end{theoremsbox}

\begin{reviewbox}
    \begin{itemize}
        \citem \textbf{Textbook Exercise Revision (Section 11.5):} Problems 5, 7, 13, 23, 37, 41.
        \begin{itemize}
            \citem Exercise 5
            \citem Exercise 7
            \citem Exercise 13
            \citem Exercise 23
            \citem Exercise 37
            \citem Exercise 41
        \end{itemize}
    \end{itemize}
\end{reviewbox}

% -------------------------------------------------------------------------
% SECTION 11.6 & 11.7: RATIO AND ROOT TESTS, TESTING STRATEGY
% -------------------------------------------------------------------------
\section{Ratio and Root Tests, Testing Strategy (Sections 11.6 \& 11.7)}

\begin{challengebox}
    \begin{itemize}
        \citem \textbf{The Ratio Test and Proof:} Write the Ratio Test ($\lim_{n \to \infty} \left|\frac{a_{n+1}}{a_n}\right| = L$: $L < 1 \implies$ absolutely convergent; $L > 1 \implies$ divergent; $L = 1 \implies$ inconclusive), and try to prove it.
    \end{itemize}
\end{challengebox}

\begin{theoremsbox}
    \begin{itemize}
        \citem \textbf{The Root Test:} Write the Root Test ($\lim_{n \to \infty} \sqrt[n]{|a_n|} = L$: $L < 1 \implies$ absolutely convergent; $L > 1 \implies$ divergent; $L = 1 \implies$ inconclusive).
    \end{itemize}
\end{theoremsbox}

\begin{warningbox}
    \begin{itemize}
        \citem \textbf{Strategy for Testing Series:} Try to write the comprehensive strategy and decision rules for testing series convergence and divergence:
        \begin{enumerate}
            \item If the $n$th term does not approach 0, use the Test for Divergence.
            \item If of the form $\sum 1/n^p$ or $\sum a r^{n-1}$, use $p$-series or geometric series test.
            \item If similar to a $p$-series or geometric series, use Direct or Limit Comparison.
            \item If the series alternates in sign, consider the Alternating Series Test or Absolute Convergence.
            \item If involving factorials or other products ($k^n, n!$), use the Ratio Test.
            \item If involving $n$th powers $(a_n = (b_n)^n)$, use the Root Test.
            \item If $f(x)$ is easily integrable, apply the Integral Test.
        \end{enumerate}
    \end{itemize}
\end{warningbox}

\begin{reviewbox}
    \begin{itemize}
        \citem \textbf{Textbook Exercise Revision (Sections 11.6 \& 11.7):}
        \begin{itemize}
            \citem \textbf{Section 11.6:} Problems 1, 5, 21, 25, 35, 41.
            \begin{itemize}
                \citem Exercise 1
                \citem Exercise 5
                \citem Exercise 21
                \citem Exercise 25
                \citem Exercise 35
                \citem Exercise 41
            \end{itemize}
            \citem \textbf{Section 11.7:} Problems 1, 3, 5, 7, 11, 13, 15, 19, 23.
            \begin{itemize}
                \citem Exercise 1
                \citem Exercise 3
                \citem Exercise 5
                \citem Exercise 7
                \citem Exercise 11
                \citem Exercise 13
                \citem Exercise 15
                \citem Exercise 19
                \citem Exercise 23
            \end{itemize}
        \end{itemize}
    \end{itemize}
\end{reviewbox}

% -------------------------------------------------------------------------
% SECTION 11.8: POWER SERIES
% -------------------------------------------------------------------------
\section{Power Series (Section 11.8)}

\begin{recallbox}
    \begin{itemize}
        \citem \textbf{Power Series in $x$:} Write the definition of a power series in $x$ ($\sum_{n=0}^\infty c_n x^n = c_0 + c_1 x + c_2 x^2 + \dots$).
    \end{itemize}
\end{recallbox}

\begin{challengebox}
    \begin{itemize}
        \citem \textbf{Convergence Interval Calculation:} For what values of $x$ is the series:
        \[ \sum_{n=0}^\infty \frac{x^n}{n!} \]
        convergent?
    \end{itemize}
\end{challengebox}

\begin{theoremsbox}
    \begin{itemize}
        \citem \textbf{Trichotomy Theorem for Power Series Convergence:} Notice the theorem that the set on which a power series $\sum_{n=0}^\infty c_n x^n$ converges is always an interval of one of the following three types:
        \begin{enumerate}
            \item[(i)] The single point $x = 0$.
            \item[(ii)] An interval $(-R, R)$, plus possibly one or both end points.
            \item[(iii)] The whole real line $(-\infty, \infty)$.
        \end{enumerate}
    \end{itemize}
\end{theoremsbox}

\begin{recallbox}
    \begin{itemize}
        \citem \textbf{Radius of Convergence:} Write the definition of the radius of convergence $R$.
        \citem \textbf{Interval of Convergence:} Write the definition of the interval of convergence of a power series.
        \citem \textbf{Power Series Centered at $a$:} Write the definition of a power series in $x$ centered at $a$:
        \[ \sum_{n=0}^\infty c_n (x - a)^n \]
    \end{itemize}
\end{recallbox}

\begin{theoremsbox}
    \begin{itemize}
        \citem \textbf{Convergence Cases for Series Centered at $a$:} Write the possible cases of interval of convergence for a power series centered at $a$: (i) single point $\{a\}$ ($R=0$); (ii) interval $(a-R, a+R)$, $[a-R, a+R)$, $(a-R, a+R]$, or $[a-R, a+R]$ ($R>0$); (iii) $(-\infty, \infty)$ ($R=\infty$).
    \end{itemize}
\end{theoremsbox}

\begin{challengebox}
    \begin{itemize}
        \citem \textbf{Radius and Interval Calculation:} Find the interval of convergence and radius of convergence of:
        \[ \sum_{n=0}^\infty \frac{(-3)^n x^n}{\sqrt{n+1}} \]
    \end{itemize}
\end{challengebox}

\begin{reviewbox}
    \begin{itemize}
        \citem \textbf{Textbook Exercise Revision (Section 11.8):} Problems 7, 23, 25, 37.
        \begin{itemize}
            \citem Exercise 7
            \citem Exercise 23
            \citem Exercise 25
            \citem Exercise 37
        \end{itemize}
    \end{itemize}
\end{reviewbox}

% -------------------------------------------------------------------------
% SECTION 11.9: REPRESENTATIONS OF FUNCTIONS AS POWER SERIES
% -------------------------------------------------------------------------
\section{Representations of Functions as Power Series (Section 11.9)}

\begin{challengebox}
    \begin{itemize}
        \citem \textbf{Power Series Representation of $\frac{1}{1+x^2}$:} Express $\frac{1}{1+x^2}$ as the sum of a power series, and also write the interval of convergence.
        \citem \textbf{Power Series Representation of $\frac{x^3}{x+2}$:} Express $\frac{x^3}{x+2}$ as the sum of a power series, and also write the interval of convergence.
    \end{itemize}
\end{challengebox}

\begin{theoremsbox}
    \begin{itemize}
        \citem \textbf{Term-by-Term Differentiation and Integration Theorem:} Notice the theorem that for a power series, after term-by-term differentiation and integration, the radius of convergence remains the same.
    \end{itemize}
\end{theoremsbox}

\begin{warningbox}
    \begin{itemize}
        \citem \textbf{Endpoint Behavior Under Differentiation/Integration:} Notice that although term-by-term differentiation and integration preserve the radius of convergence $R$, the interval of convergence may change at the endpoints.
    \end{itemize}
\end{warningbox}

\begin{challengebox}
    \begin{itemize}
        \citem \textbf{Interval and Explicit Sum Formula:} Determine the interval of convergence for:
        \[ x + \frac{x^2}{2} + \frac{x^3}{3} + \frac{x^4}{4} + \dots \]
        and find an explicit formula for this series when it is convergent.
        \citem \textbf{Power Series for $\tan^{-1}x$:} Find the power series representation for $\tan^{-1}x$, and also find its interval of convergence.
        \citem \textbf{Definite Integral Approximation:} Approximate:
        \[ \int_0^{1/2} \frac{1}{1+x^7} \, dx \]
        correct to within $10^{-7}$.
    \end{itemize}
\end{challengebox}

\begin{reviewbox}
    \begin{itemize}
        \citem \textbf{Textbook Exercise Revision (Section 11.9):} Problems 5, 7, 9, 21, 27, 31.
        \begin{itemize}
            \citem Exercise 5
            \citem Exercise 7
            \citem Exercise 9
            \citem Exercise 21
            \citem Exercise 27
            \citem Exercise 31
        \end{itemize}
    \end{itemize}
\end{reviewbox}

% -------------------------------------------------------------------------
% SECTION 11.10: TAYLOR AND MACLAURIN SERIES
% -------------------------------------------------------------------------
\section{Taylor and Maclaurin Series (Section 11.10)}

\begin{challengebox}
    \begin{itemize}
        \citem \textbf{Taylor Coefficient Formula Proof:} Prove that if $f$ has a power series representation (expansion) at $a$:
        \[ f(x) = c_0 + c_1(x - a) + c_2(x - a)^2 + c_3(x - a)^3 + \dots \quad \text{when } |x - a| < R, \]
        then:
        \[ c_n = \frac{f^{(n)}(a)}{n!} \]
    \end{itemize}
\end{challengebox}

\begin{recallbox}
    \begin{itemize}
        \citem \textbf{Taylor and Maclaurin Series Definitions:} Write the definition of the Taylor series of $f$ at $a$ ($\sum_{n=0}^\infty \frac{f^{(n)}(a)}{n!}(x-a)^n$), and also the definition of the Maclaurin series ($a = 0$).
        \citem \textbf{$n$th-Degree Taylor Polynomial:} Write the definition of the $n$th-degree Taylor polynomial $T_n(x)$ of $f$ at $a$.
        \citem \textbf{Sum of Taylor Series Condition:} When is $f(x)$ called the sum of its Taylor series ($f(x) = \lim_{n \to \infty} T_n(x)$)?
        \citem \textbf{Remainder of Taylor Series:} Write the definition of the remainder $R_n(x) = f(x) - T_n(x)$ of a Taylor series.
    \end{itemize}
\end{recallbox}

\begin{theoremsbox}
    \begin{itemize}
        \citem \textbf{Taylor Series Convergence Theorem:} Notice the theorem that if $f(x) = T_n(x) + R_n(x)$ and:
        \[ \lim_{n \to \infty} R_n(x) = 0 \]
        for $|x - a| < R$, then $f$ is equal to the sum of its Taylor series on the interval $|x - a| < R$.
        \citem \textbf{Taylor's Inequality:} In order to show $\lim_{n \to \infty} R_n(x) = 0$, we can use Taylor's Inequality. Please state what it is: if $|f^{(n+1)}(x)| \le M$ for $|x - a| \le d$, then the remainder $R_n(x)$ of the Taylor series satisfies:
        \[ |R_n(x)| \le \frac{M}{(n+1)!} |x - a|^{n+1} \quad \text{for } |x - a| \le d. \]
        \citem \textbf{Factorial Limit Identity:} Notice that:
        \[ \lim_{n \to \infty} \frac{x^n}{n!} = 0 \quad \text{for all real } x. \]
    \end{itemize}
\end{theoremsbox}

\begin{challengebox}
    \begin{itemize}
        \citem \textbf{Maclaurin Series for $e^x$ Proof:} Prove that $e^x$ is equal to the sum of its Maclaurin series for all real $x$ by using Taylor's Inequality.
        \citem \textbf{Binomial Series Expansion and Convergence:} Find the Maclaurin series for $(1+x)^k$ and show that it is convergent for $|x| < 1$.
        \citem \textbf{Maclaurin Series of $\frac{1}{\sqrt{4 - x}}$:} Find the Maclaurin series of:
        \[ \frac{1}{\sqrt{4 - x}} \]
        and determine its radius of convergence.
    \end{itemize}
\end{challengebox}

\begin{warningbox}
    \begin{itemize}
        \citem \textbf{Important Maclaurin Series Reference:} Refer back to key Maclaurin series in lecture notes ($e^x, \sin x, \cos x, \frac{1}{1-x}, \ln(1+x), \tan^{-1}x, (1+x)^k$).
    \end{itemize}
\end{warningbox}

\begin{challengebox}
    \begin{itemize}
        \citem \textbf{Definite Integral Power Series Approximation:} Evaluate:
        \[ \int_0^1 e^{-x^2} \, dx \]
        correct to within an error of $0.001$.
    \end{itemize}
\end{challengebox}

\begin{reviewbox}
    \begin{itemize}
        \citem \textbf{Textbook Exercise Revision (Section 11.10):} Problems 15, 21, 37, 45, 57, 61, 63, 67, 79, 81, 83, 85.
        \begin{itemize}
            \citem Exercise 15
            \citem Exercise 21
            \citem Exercise 37
            \citem Exercise 45
            \citem Exercise 57
            \citem Exercise 61
            \citem Exercise 63
            \citem Exercise 67
            \citem Exercise 79
            \citem Exercise 81
            \citem Exercise 83
            \citem Exercise 85
        \end{itemize}
    \end{itemize}
\end{reviewbox}

% -------------------------------------------------------------------------
% SECTION 11.11: APPLICATIONS OF TAYLOR POLYNOMIALS
% -------------------------------------------------------------------------
\section{Applications of Taylor Polynomials (Section 11.11)}

\begin{challengebox}
    \begin{itemize}
        \citem \textbf{Taylor Polynomial Approximation and Accuracy Bound:} Approximate the function $f(x) = \sqrt[3]{x}$ by a Taylor polynomial of degree 2 at $a = 8$, and determine how accurate this approximation is when $7 \le x \le 9$.
    \end{itemize}
\end{challengebox}

\begin{warningbox}
    \begin{itemize}
        \citem \textbf{Alternating Series Error Bound Note:} For an alternating series where the terms $b_n$ are not necessarily positive, notice that:
        \[ |S - S_n| \le |b_{n+1}| \]
    \end{itemize}
\end{warningbox}

\begin{reviewbox}
    \begin{itemize}
        \citem \textbf{Textbook Exercise Revision (Section 11.11):} Problems 5, 13, 19, 25, 28.
        \begin{itemize}
            \citem Exercise 5
            \citem Exercise 13
            \citem Exercise 19
            \citem Exercise 25
            \citem Exercise 28
        \end{itemize}
    \end{itemize}
\end{reviewbox}

% =========================================================================
% PART 2: CHAPTER 10 — PARAMETRIC EQUATIONS AND POLAR COORDINATES
% =========================================================================

% -------------------------------------------------------------------------
% SECTION 10.1: CURVES DEFINED BY PARAMETRIC EQUATIONS
% -------------------------------------------------------------------------
\section{Curves Defined by Parametric Equations (Section 10.1)}

\begin{recallbox}
    \begin{itemize}
        \citem \textbf{Parametric Curve Definitions:} We can define a curve by parametric equations $x = f(t), y = g(t)$ ($a \le t \le b$). What is the initial point, and what is the terminal point?
    \end{itemize}
\end{recallbox}

\begin{challengebox}
    \begin{itemize}
        \citem \textbf{Cycloid and Parametric Derivation:} What is a cycloid, and try to derive the parametric equations of this curve ($x = r(\theta - \sin\theta), y = r(1 - \cos\theta)$).
    \end{itemize}
\end{challengebox}

\begin{reviewbox}
    \begin{itemize}
        \citem \textbf{Textbook Exercise Revision (Section 10.1):} Problems 17, 37, 47, 51.
        \begin{itemize}
            \citem Exercise 17
            \citem Exercise 37
            \citem Exercise 47
            \citem Exercise 51
        \end{itemize}
    \end{itemize}
\end{reviewbox}

% -------------------------------------------------------------------------
% SECTION 10.2: CALCULUS WITH PARAMETRIC CURVES
% -------------------------------------------------------------------------
\section{Calculus with Parametric Curves (Section 10.2)}

\begin{theoremsbox}
    \begin{itemize}
        \citem \textbf{First Derivative of Parametric Equations:} For parametric equations $x = f(t), y = g(t)$, how do we define:
        \[ \frac{dy}{dx} = \frac{dy/dt}{dx/dt} = \frac{g'(t)}{f'(t)} \quad \left(\text{provided } \frac{dx}{dt} \neq 0\right) \]
        \citem \textbf{Second Derivative of Parametric Curves:} How do we find:
        \[ \frac{d^2y}{dx^2} = \frac{\od{}{t}\left(\frac{dy}{dx}\right)}{\frac{dx}{dt}} \]
    \end{itemize}
\end{theoremsbox}

\begin{warningbox}
    \begin{itemize}
        \citem \textbf{Parametric Second Derivative Pitfall:} Notice that $\frac{d^2y}{dx^2} \neq \frac{d^2y/dt^2}{d^2x/dt^2}$. Always differentiate $\frac{dy}{dx}$ with respect to $t$ and divide by $\frac{dx}{dt}$.
    \end{itemize}
\end{warningbox}

\begin{theoremsbox}
    \begin{itemize}
        \citem \textbf{Area Under a Parametric Curve:} Write the formula for the area under the parametric curve $x = f(t), y = g(t), a \le t \le b$, where $g(t) \ge 0$:
        \[ A = \int_a^b y \, dx = \int_\alpha^\beta g(t) f'(t) \, dt \]
        \citem \textbf{Arc Length of a Parametric Curve:} Write the arc length formula for the parametric curve $x = f(t), y = g(t), a \le t \le b$:
        \[ L = \int_a^b \sqrt{\left(\frac{dx}{dt}\right)^2 + \left(\frac{dy}{dt}\right)^2} \, dt \]
    \end{itemize}
\end{theoremsbox}

\begin{reviewbox}
    \begin{itemize}
        \citem \textbf{Textbook Exercise Revision (Section 10.2):} Problems 5, 7, 19, 23, 29, 35, 37, 39, 47.
        \begin{itemize}
            \citem Exercise 5
            \citem Exercise 7
            \citem Exercise 19
            \citem Exercise 23
            \citem Exercise 29
            \citem Exercise 35
            \citem Exercise 37
            \citem Exercise 39
            \citem Exercise 47
        \end{itemize}
    \end{itemize}
\end{reviewbox}

% -------------------------------------------------------------------------
% SECTION 10.3: POLAR COORDINATES
% -------------------------------------------------------------------------
\section{Polar Coordinates (Section 10.3)}

\begin{recallbox}
    \begin{itemize}
        \citem \textbf{Polar Coordinates Definition:} Write the definition of polar coordinates $(r, \theta)$, the pole, and the polar axis, along with the standard conversions $x = r\cos\theta, y = r\sin\theta, r^2 = x^2 + y^2, \tan\theta = y/x$.
        \citem \textbf{Lima\c{c}ons and Cardioid Expressions:} Write the polar graph expressions of lima\c{c}ons ($r = a \pm b\cos\theta, r = a \pm b\sin\theta$) and the cardioid ($a = b$).
        \citem \textbf{Rose Curves and Petal Rules:} Write the expressions for rose curves ($r = a\cos(n\theta), r = a\sin(n\theta)$), and state how to determine how many leaves (petals) these roses have ($n$ petals if $n$ is odd, $2n$ petals if $n$ is even).
    \end{itemize}
\end{recallbox}

\begin{reviewbox}
    \begin{itemize}
        \citem \textbf{Textbook Exercise Revision (Section 10.3):} Problems 3, 5, 7, 9.
        \begin{itemize}
            \citem Exercise 3
            \citem Exercise 5
            \citem Exercise 7
            \citem Exercise 9
        \end{itemize}
    \end{itemize}
\end{reviewbox}

% -------------------------------------------------------------------------
% SECTION 10.4: AREAS AND LENGTHS IN POLAR COORDINATES
% -------------------------------------------------------------------------
\section{Areas and Lengths in Polar Coordinates (Section 10.4)}

\begin{theoremsbox}
    \begin{itemize}
        \citem \textbf{Polar Region Area Formula:} How do we find the area of a region in polar coordinates? Write the formula:
        \[ A = \int_a^b \frac{1}{2} r^2 \, d\theta = \int_a^b \frac{1}{2} [f(\theta)]^2 \, d\theta \]
        \citem \textbf{Polar Curve Arc Length Formula:} How do we find the arc length of a polar graph in polar coordinates? Write the formula:
        \[ L = \int_a^b \sqrt{r^2 + \left(\frac{dr}{d\theta}\right)^2} \, d\theta \]
        \citem \textbf{Polar Graph Tangent Line Formula:} How do we find the tangent line of a polar graph in polar coordinates? Write the slope formula:
        \[ \frac{dy}{dx} = \frac{\frac{dr}{d\theta}\sin\theta + r\cos\theta}{\frac{dr}{d\theta}\cos\theta - r\sin\theta} \]
    \end{itemize}
\end{theoremsbox}

\begin{reviewbox}
    \begin{itemize}
        \citem \textbf{Textbook Exercise Revision (Section 10.4):} Problems 1, 3, 5, 17, 23, 27, 31, 49, 57, 71.
        \begin{itemize}
            \citem Exercise 1
            \citem Exercise 3
            \citem Exercise 5
            \citem Exercise 17
            \citem Exercise 23
            \citem Exercise 27
            \citem Exercise 31
            \citem Exercise 49
            \citem Exercise 57
            \citem Exercise 71
        \end{itemize}
    \end{itemize}
\end{reviewbox}

\end{document}
